\uuid{0s2X}
\titre{Image affine d'une variable normale}
\auteur{Maxime Nguyen}
\organisation{AMSCC}
\datecreate{2026-11-27}
\niveau{L2}
\module{Probabilité}
\chapitre{Probabilité continue}
\sousChapitre{Loi normale}
\difficulte{1}
\theme{loi normale, transformation affine, écart-type}

\contenu{
\begin{question}
\qtype{onlyone}
On note $\mathcal{N}(\mu,\sigma)$ la loi normale d'espérance $\mu$ et d'écart-type $\sigma > 0$. Soit $X \sim \mathcal{N}(2, 3)$ et $Y = -2X + 1$. La loi de $Y$ est :
\begin{answers}
\bad{$\mathcal{N}(-3, -6)$}
\bad{$\mathcal{N}(5, 6)$}
\good{$\mathcal{N}(-3, 6)$}
\bad{$\mathcal{N}(-3, 36)$}
\bad{$\mathcal{N}(-3, 7)$}
\end{answers}
\begin{explanations}
Une image affine $aX + b$ d'une variable normale
($a \neq 0$) est normale, d'espérance $a\mu + b = -4 + 1 = -3$ et d'écart-type
$|a|\,\sigma = 2 \times 3 = 6$.

A : un écart-type est positif ; c'est $|a|$ qui intervient.

B : erreur de signe sur l'espérance ($2 \times 2 + 1$).

D donne la variance $36$ à la place de l'écart-type.

E ajoute la constante $1$ à l'écart-type ; une translation ne change pas la dispersion.

Contre-sens visés : écart-type négatif (A) ; confusion écart-type /
variance (D) ; croire que la constante modifie la dispersion (E).
\end{explanations}
\end{question}
}
